\documentclass[11pt,twocolumn]{article}
\usepackage[letterpaper,margin=0.7in,columnsep=0.25in]{geometry}
\usepackage[T1]{fontenc}
\usepackage{times,microtype,enumitem,xcolor,hyperref}
\hypersetup{colorlinks=true,urlcolor=blue!55!black,citecolor=blue!55!black,linkcolor=black,
pdftitle={What Does the Image Add? Learning to Supplement Existing Descriptions},pdfauthor={}}
\setlength{\parindent}{0.8em}
\setlength{\parskip}{0pt}
\clubpenalty=10000
\widowpenalty=10000
\setlist{nosep,leftmargin=*}
\makeatletter
\renewcommand\section{\@startsection{section}{1}{0pt}{8pt plus 2pt minus 1pt}{4pt}{\large\bfseries}}
\makeatother
\begin{document}
\twocolumn[
\begin{center}
{\Large\bfseries What Does the Image Add?\\[2pt]
Learning to Supplement Existing Descriptions\par}
\vspace{5pt}
{\small CSCI-GA 2271 Computer Vision \quad Fall 2026 \quad Project Proposal}\par
\vspace{3pt}
{\small\bfseries Track: Agent Track \qquad Format: Research Project}\par
\vspace{3pt}
{\small September 29, 2026}\par
\end{center}
\vspace{3pt}
]

\section{Project Idea and Motivation}
A short image description often leaves out details that a reader might find
useful. We want to study whether a vision-language model can look at an image
and an existing description, then add a small amount of \textbf{correct, useful
information that the text does not already contain}.

For example, suppose a photograph shows a red, long-sleeved shirt. If the text
says ``a shirt,'' a useful addition could mention its color and sleeves. If the
text already says ``a red shirt,'' the model should focus on the sleeves.
The image is unchanged, but what is worth saying depends on what the reader
already knows.

This could support interactive image exploration or help enrich short image
descriptions. We will investigate whether the added details actually help
readers answer questions about an image. Our main research question is:
\textbf{can a model learn to supplement existing text with visually supported
facts, while avoiding repetition and unsupported claims?}

\section{Background and Research Focus}
Related work provides several starting points. FINECAPTION studies descriptions
of selected image regions \cite{finecaption}. CapWAP trains captions to serve
information needs \cite{capwap}, and PromptCap uses a question to guide what a
caption describes \cite{promptcap}.

We will focus on the information already present in a description. Our
hypothesis is that training with different descriptions of the same image can
help a model learn \emph{what remains to be said}. At generation time, the model
will not see the downstream evaluation question. We will also ask whether
pointing the model toward an object helps, or whether looking directly at the
whole image is sufficient. The intended contribution is a method and empirical
study of this text-dependent information gap; its novelty and usefulness will
need to be established through comparison with existing approaches.

\section{Proposed Method}
\textbf{Starting system.} We will use an open pretrained vision-language model,
such as Qwen2.5-VL-3B, to take an image and an existing description and generate
a concise supplement. We will begin with objects mentioned in the description
to keep the initial task well defined. LoRA fine-tuning will allow us to adapt
the model without training it from scratch.

\textbf{Learning what is missing.} We plan to construct training examples in
which the same object appears with different amounts of information in the
text. An object with several annotated attributes can provide multiple
examples: when the text mentions one attribute, the target supplement can
contain another. Natural captions describing the same object will provide a
second source of examples. We will check that a target adds meaning, since
different wording can express the same fact.

The initial training objective will teach the model to generate these missing
details. We will investigate paired examples where a fact changes from absent
to already mentioned, encouraging the model to adjust its output accordingly.
We will compare this with ordinary description fine-tuning to find out whether
explicitly modeling the existing text provides a benefit.

\textbf{Whole image or region guidance?} Our main comparison will give one
version the full image and text, and another the same inputs with a box marking
the object of interest. Both will receive the same textual object reference and
comparable training data, initialization and compute. This will let us test
whether boxes help rather than assume they are necessary. If they help, we will
explore automatic localization; if they do not, we will favor the simpler
full-image approach. Letting the model choose which object to discuss will be
an optional extension after the initial task is understood.

\section{Data and Evaluation}
\textbf{Datasets.} We plan to use Visual Genome's object and attribute
annotations \cite{krishna}, GQA's associated questions \cite{hudson}, and
Flickr30k Entities' captions and object correspondences \cite{plummer}.
Controlled attribute examples will help us examine model behavior, while
natural captions will test whether it works with ordinary descriptions.
Visual Genome and GQA share imagery, so we will keep their splits aligned.
Training and evaluation will be separated by image, with duplicate checks and
a reserved set for final evaluation.

\textbf{Comparisons.} Baselines will include the existing text alone, a generic
image caption, a model prompted to add details from the full image, and ordinary
fine-tuning. We will compare versions with and without region guidance under
matched training conditions and output lengths. A more detailed full-image
caption will also test whether simply saying more is enough.

\textbf{What counts as success?} We will evaluate:
\begin{itemize}
\item \textbf{New information:} how many reference facts missing from the
original text are recovered, and how much existing information is repeated.
\item \textbf{Correctness:} whether additions are supported by the image and
refer to the right object, including contradictions and other factual errors.
\item \textbf{Usefulness:} whether a text-only question-answering model answers
independent image questions more accurately after receiving the supplement.
\item \textbf{Cost:} generation time, model calls, output length, memory use
and total GPU-hours, including any localization step.
\end{itemize}

Questions and answers will be hidden from supplement generation. We will
report cases where an addition makes an answer worse as well as cases where it
helps. Reference matching and model-based judgments will support evaluation,
but neither alone will establish factual correctness. Automated reviewers will
be checked on clearly correct and incorrect examples, and a visually reviewed
sample will examine their reliability. We will report uncertainty, repeat key
training runs, and test transfer across the two image sources.

\section{Expected Deliverables and Risks}
We expect to produce a reproducible training and evaluation pipeline, a
comparison of the proposed method and full-image baselines, and an analysis of
when region guidance helps. The final submission will include a 4--9-page
LaTeX research report excluding references, a code repository with incremental
history, presentation materials, and an Agent Log appendix. A small viewer for
comparing descriptions will be a stretch goal.

The main risks are incomplete annotations, repetitive or invented details,
and a gap between controlled examples and natural language. We will address
these through data inspection, factual-error analysis, and evaluation on
natural descriptions. A simple full-image baseline may already perform well;
if so, understanding its strengths and the limits of additional training will
be part of the research outcome. We will use development results to revise
unsuccessful ideas and keep the final evaluation separate from method selection.

\section{Timeline, Resources, and Team Workflow}
\textbf{Preliminary timeline} (aligned to course deadlines):
\begin{itemize}
\item \textbf{Weeks 1--2:} Review related work, prepare image-level splits,
inspect training examples, and implement baseline models.
\item \textbf{Weeks 3--4:} Train the supplement model and compare ordinary
training, paired examples, and versions with and without boxes.
\item \textbf{Weeks 5--6:} Evaluate natural descriptions, question-answering
utility, factual errors and transfer; investigate the main failure cases.
\item \textbf{Final two weeks:} Complete final evaluation, report and Agent
Log; prepare the presentation and team code walkthrough.
\end{itemize}

\textbf{Resources.} We plan to use NYU HPC with one suitable available GPU at a
time and a small capped pilot before scaling up. We will adapt batching and
image resolution to available memory, cache reusable work, and monitor runtime
and utilization to avoid wasting allocated compute.

\textbf{Agent Track workflow.} Agents will assist with literature review,
implementation and analysis. Team members will review research decisions,
inspect failures and verify the final system. The Agent Log will document the
tools used, division of work, agent mistakes, human diagnosis and verification.
Individual contributions will be listed in the final report, and every member
will prepare to explain the submitted code.

\begin{thebibliography}{6}
\small
\bibitem{finecaption} H. Hua et al. FINECAPTION: Compositional Image Captioning
Focusing on Wherever You Want at Any Granularity. 2024.
\href{https://arxiv.org/abs/2411.15411}{arXiv:2411.15411}.
\bibitem{capwap} A. Fisch et al. CapWAP: Image Captioning with a Purpose.
\emph{EMNLP}, 2020.
\href{https://aclanthology.org/2020.emnlp-main.705/}{ACL Anthology}.
\bibitem{promptcap} Y. Hu et al. PromptCap: Prompt-Guided Task-Aware Image
Captioning. \emph{ICCV}, 2023.
\href{https://arxiv.org/abs/2211.09699}{arXiv:2211.09699}.
\bibitem{krishna} R. Krishna et al. Visual Genome: Connecting Language and Vision
Using Crowdsourced Dense Image Annotations. 2016.
\href{https://arxiv.org/abs/1602.07332}{arXiv:1602.07332}.
\bibitem{hudson} D. A. Hudson and C. D. Manning. GQA: A New Dataset for Real-World
Visual Reasoning and Compositional Question Answering. \emph{CVPR}, 2019.
\href{https://arxiv.org/abs/1902.09506}{arXiv:1902.09506}.
\bibitem{plummer} B. A. Plummer et al. Flickr30k Entities: Collecting Region-to-Phrase
Correspondences for Richer Image-to-Sentence Models. 2015.
\href{https://arxiv.org/abs/1505.04870}{arXiv:1505.04870}.
\end{thebibliography}
\end{document}
